\section{Optimality and the equality caps}\label{sec:equality}

The cap theorem now applies separately to the two caps of the
original body. Its equality statement also forces both effective
angles to be full quarter turns. The following section uses these caps to recover the hull and the
body itself, including every boundary point.

\begin{theorem}[Optimality]\label{thm:optimality}
Every ambidextrous sofa has area at most $\Mc$. The value $\Mc$ is
attained by $\SR$.
\end{theorem}

\begin{proof}
The small-area branch of \cref{prop:reduction} is strictly below
$\Mc$. This comparison uses exact inequalities:
\[
 \sqrt2<\frac85<\Mc.
\]
For example, $Y>2/7$, the cubic equation, and
$\arctan Y\ge Y-Y^3/3$ give
\[
 \Mc\ge\frac{11}{12}+4Y^2+\frac54Y
       >\frac{11}{12}+\frac{16}{49}+\frac5{14}
       =\frac{941}{588}>\frac85.
\]
The arctangent estimate follows by differentiating
$\arctan y-y+y^3/3$ on $[0,\infty)$; its derivative is
$y^4/(1+y^2)\ge0$.

In the remaining branch, the reduction gives
\[
 |S|\le P_\alpha(U)+P_\gamma(V).
\]
Both caps satisfy the hypotheses of \cref{thm:cap-bound}, so each
score is at most $\Mc/2$. This proves the upper bound. The explicit
construction in \cref{thm:reference} supplies an ambidextrous sofa
$\SR$ of area $\Mc$.
\end{proof}

\begin{proposition}[The two original equality caps]\label{prop:equality-caps}
If an ambidextrous sofa has area $\Mc$, then in the normalization
of \cref{sec:reduction} its two caps satisfy
\[
 \alpha=\gamma=\frac\pi2,\qquad
 U=\KR+(a,0),\qquad V=\KR+(b,0)
 \quad\text{for some }a,b\in\R.
\]
All four terms in the deficit identity \eqref{eq:four-term-deficit}
vanish for $c=\Mc$.
\end{proposition}
\begin{proof}

Suppose $|S|=\Mc$. In the normalization of \cref{sec:reduction},
$|T|=\Mc>\sqrt2$. Apply the exact deficit
\eqref{eq:four-term-deficit} with $c=\Mc$. Every term on its
right-hand side is nonnegative: two by the scalar cap theorem,
one by the positive-part formula \eqref{eq:partition-defect}, and
one because it is the area of a set. Hence all four vanish:
\begin{equation}\label{eq:equality-deficits}
 P_\alpha(U)=P_\gamma(V)=\frac\Mc2,\qquad
 D_{\rm part}=0,\qquad |E\setminus T|=0.
\end{equation}
In particular, \cref{thm:cap-equality} gives
\begin{equation}\label{eq:original-equality-caps}
 \alpha=\gamma=\frac\pi2,\qquad
 U=\KR+(a,0),\qquad V=\KR+(b,0)
 \quad\text{for some }a,b\in\R.
\end{equation}
These are the caps defined from $\conv T$ in
\eqref{eq:actual-caps}.

For completeness, vanishing of the partition defect gives the
following pointwise relations almost everywhere:
\[
 \begin{array}{ll}
 d_V\le N_\alpha(U),\quad d_U\le N_\gamma(V),& x\in J,\\
 N_\alpha(U)\le d_V,\quad N_\gamma(V)\le d_U,&x\in I\setminus J.
 \end{array}
\]
The vanishing filling defect says that the physical envelope differs
from $T$ by a null set. Exact recovery requires the additional
topological argument below.


\end{proof}
